\section{Methodology}

\subsection{Algorithmic Redistricting}

We use recursive bisection via METIS graph partitioning~\cite{karypis1998metis,deluca2026recursive} to generate algorithmic redistricting plans for all 50 states. The algorithm:
\begin{itemize}
\item Represents census tracts as graph vertices
\item Connects adjacent tracts with edges
\item Recursively bisects the graph to create balanced districts
\item Optimizes for edge-cut minimization (compactness proxy)
\item Cannot access partisan data during partitioning
\end{itemize}

\subsection{Efficiency Gap Calculation}

For each district $i$ with Democratic votes $D_i$ and Republican votes $R_i$:

\[
\text{Wasted}_D^i = \begin{cases}
D_i - \frac{D_i + R_i}{2} - 1 & \text{if Democrat wins} \\
D_i & \text{if Republican wins}
\end{cases}
\]

\[
\text{Wasted}_R^i = \begin{cases}
R_i & \text{if Democrat wins} \\
R_i - \frac{D_i + R_i}{2} - 1 & \text{if Republican wins}
\end{cases}
\]

State-level efficiency gap:
\[
\text{EG} = \frac{\sum_i \text{Wasted}_R^i - \sum_i \text{Wasted}_D^i}{\sum_i (D_i + R_i)} \times 100\%
\]

Positive values favor Republicans; negative values favor Democrats.

\subsection{Data}

\begin{itemize}
\item \textbf{Districts}: Algorithmic plans from recursive bisection; enacted plans from state legislatures (2020 cycle)
\item \textbf{Elections}: Presidential (2016, 2020), midterm (2018) results at precinct level
\item \textbf{States}: All 50 states; focus on 15 competitive states for detailed analysis
\end{itemize}

\subsection{Comparison Framework}

We compute efficiency gaps for algorithmic and enacted plans using identical election data, enabling direct comparison of partisan bias while holding geography constant.
